\chapter{Modelado y diseño de soluciones algorítmicas}
\label{cap:02-diseno-algoritmico}
\chaptermark{Diseño algorítmico}

Una vez comprendido y formulado un problema, es necesario establecer un procedimiento que transforme los datos disponibles en los resultados requeridos. El diseño algorítmico constituye el vínculo entre el análisis realizado en el capítulo anterior y la implementación que se abordará en el capítulo siguiente. Su propósito es construir una solución precisa, comprensible y susceptible de comprobación antes de escribir código en un lenguaje de programación. \citet{zingaro2021} se propone como lectura complementaria sobre el aprendizaje de la programación mediante problemas.

\citet[pp.~9–10]{farrell2013} sitúa la planificación de la lógica antes de la codificación y destaca el uso del pseudocódigo, los diagramas de flujo y las pruebas de escritorio. En este capítulo se adopta ese enfoque mediante ejemplos originales vinculados con situaciones académicas, comerciales y cotidianas.

\begin{resultados}
  \begin{itemize}
    \item Traducir requerimientos a un modelo de entrada, proceso y salida.
    \item Diseñar algoritmos mediante pseudocódigo y diagramas de flujo.
    \item Comprobar decisiones y ciclos con pruebas de escritorio.
  \end{itemize}
\end{resultados}

\section{Concepto y características de los algoritmos}

Un algoritmo es una secuencia organizada de instrucciones que permite resolver un problema o realizar una tarea definida. En términos computacionales, especifica qué operaciones deben efectuarse, en qué orden y bajo qué condiciones. Un mismo problema puede admitir varios algoritmos; la elección depende de las restricciones, la claridad y el costo de ejecución.

Por ejemplo, para calcular el importe de una compra se necesita conocer el precio unitario y la cantidad de unidades. El algoritmo multiplica ambos valores y presenta el resultado. La multiplicación es una operación; el algoritmo comprende además la obtención de los datos, las condiciones de validez y la comunicación del importe.

\begin{tablalibro}{2}{L{0.31}L{0.69}}
\caption{Características de un algoritmo}\label{tab:caracteristicas-de-un-algoritmo}\\
\toprule
\rowcolor{Papel}
\CabeceraTabla{Característica} & \CabeceraTabla{Interpretación} \\
\midrule
\endfirsthead
\SigueTabla{2}
\toprule
\rowcolor{Papel}
\CabeceraTabla{Característica} & \CabeceraTabla{Interpretación} \\
\midrule
\endhead
\midrule
\endfoot
\bottomrule
\endlastfoot

Precisión & Cada instrucción indica una acción sin ambigüedad. \\
Finitud & Para las entradas admitidas, el procedimiento termina. \\
Efectividad & Las operaciones son realizables en un número finito de pasos. \\
Orden lógico & Las instrucciones respetan sus dependencias. \\
Entradas y salidas & Se especifican los datos requeridos, si los hay, y los resultados. \\
Corrección & El resultado satisface la especificación para toda entrada válida. \\
\end{tablalibro}


\subsection{Algoritmo, programa y proceso}

Un algoritmo es la descripción de la solución; un programa es una implementación expresada en un lenguaje ejecutable; un proceso es una instancia de ejecución de un programa. Esta distinción permite revisar la lógica sin depender de un lenguaje específico. La expresión de una solución mediante pseudocódigo tampoco garantiza que pueda ejecutarse directamente.

\subsection{Ejemplo: conversión de minutos a horas y minutos}

Si una actividad dura 145 minutos, se obtiene el cociente entero al dividir entre 60 y el residuo de esa división. El resultado es 2 horas y 25 minutos. La entrada debe ser un entero no negativo. El procedimiento es aplicable a cualquier duración que cumpla esta condición, incluso cero.

\Pseudocodigo{02-conversion-minutos.txt}

DIV representa el cociente entero y MOD el residuo. El ejemplo presupone una entrada válida; si esta no está garantizada, se requiere una comprobación previa. Esta separación entre supuestos y acciones será importante en todo el capítulo.

\section{Análisis de requerimientos y condiciones del problema}

Un requerimiento describe una necesidad que debe satisfacer la solución. Antes de diseñar instrucciones es necesario esclarecer qué se espera del sistema, cuáles son las reglas aplicables y cómo se determinará que el resultado es aceptable. Una solicitud como «calcular las notas» es insuficiente mientras no se definan la escala, las ponderaciones y la condición de aprobación.

\citet[pp.~8–10]{farrell2013} señala la importancia de comprender las necesidades del usuario y planificar los pasos de la solución. En la práctica, esa comprensión puede documentarse mediante entrevistas, ejemplos de resultados y especificaciones de los datos.

\subsection{Requerimientos funcionales, restricciones y criterios de aceptación}

Los requerimientos funcionales describen las operaciones que debe realizar la solución. Las restricciones delimitan los datos y las acciones permitidas; los criterios de aceptación establecen resultados observables que permiten comprobar el cumplimiento. Las condiciones no funcionales, como tiempo de respuesta o facilidad de uso, también pueden afectar el diseño.

\begin{tablalibro}{2}{L{0.31}L{0.69}}
\caption{Especificación de un cálculo de calificaciones}\label{tab:especificacion-de-un-calculo-de-calificaciones}\\
\toprule
\rowcolor{Papel}
\CabeceraTabla{Elemento} & \CabeceraTabla{Descripción} \\
\midrule
\endfirsthead
\SigueTabla{2}
\toprule
\rowcolor{Papel}
\CabeceraTabla{Elemento} & \CabeceraTabla{Descripción} \\
\midrule
\endhead
\midrule
\endfoot
\bottomrule
\multicolumn{2}{@{}p{\linewidth}@{}}{\NotaLibro{Elaboración propia.}}\\
\endlastfoot

Objetivo & Calcular el promedio de tres evaluaciones. \\
RF-01 & Recibir tres notas numéricas. \\
RF-02 & Calcular la media aritmética e informar la condición. \\
Restricción & Cada nota pertenece al intervalo de 0 a 20. \\
Regla & Aprobado si el promedio es mayor o igual a 11; regla didáctica del ejemplo. \\
Criterio de aceptación & Con 10, 11 y 12, el promedio es 11 y la condición es aprobado. \\
Entrada inválida & Informar el error y terminar sin calcular el promedio. \\
\end{tablalibro}


La regla de aprobación del ejemplo es una decisión explícita de modelado, no una afirmación sobre un reglamento institucional. El promedio se compara sin redondear; si se presenta con dos decimales, ese formato de salida no modifica la decisión.

\subsection{Precondiciones y poscondiciones}

Una precondición establece qué debe cumplirse antes de ejecutar una operación; una poscondición describe qué se garantiza al finalizarla correctamente. Para calcular una media de tres notas, una precondición es que las tres sean numéricas y válidas. La poscondición es que el resultado sea la suma dividida entre tres y pertenezca al intervalo de 0 a 20. Si no se cumple la precondición, la solución debe tratar el error de acuerdo con su especificación.

\section{Identificación de datos de entrada, procesamiento y resultados esperados}

El modelo entrada–proceso–salida (EPS) representa la transformación que realiza una solución. \citet[p.~9]{farrell2013} presenta una herramienta equivalente, denominada gráfica IPO por sus iniciales en inglés. Identificar estos componentes permite evitar cálculos innecesarios y detectar datos faltantes. La \RefFigura{fig:modelo-eps} resume los tres componentes.

\FiguraLibro{04-modelo-eps}{Modelo de entrada, procesamiento y salida}{fig:modelo-eps}{Elaboración propia.}

Las entradas pueden ser valores proporcionados por una persona, lecturas de sensores o registros previamente almacenados. El procesamiento incluye operaciones aritméticas, comparaciones, conteos y transformaciones. Las salidas pueden ser números, mensajes, listas o decisiones. Un valor intermedio no es necesariamente una salida: solo lo será si la especificación exige comunicarlo.

\begin{tablalibro}{3}{L{0.31}L{0.39}L{0.30}}
\caption{Modelo EPS para un pedido de cafetería}\label{tab:modelo-eps-para-un-pedido-de-cafeteria}\\
\toprule
\rowcolor{Papel}
\CabeceraTabla{Entrada} & \CabeceraTabla{Proceso} & \CabeceraTabla{Salida} \\
\midrule
\endfirsthead
\SigueTabla{3}
\toprule
\rowcolor{Papel}
\CabeceraTabla{Entrada} & \CabeceraTabla{Proceso} & \CabeceraTabla{Salida} \\
\midrule
\endhead
\midrule
\endfoot
\bottomrule
\endlastfoot

Cantidad entera positiva; precio unitario no negativo & Validar los datos y multiplicar cantidad por precio & Subtotal en soles \\
Subtotal y tasa de descuento definida & Aplicar descuento y calcular importe final & Descuento e importe final \\
Dato fuera del dominio & Interrumpir el cálculo & Mensaje de error \\
\end{tablalibro}


\subsection{Diccionario de datos}

El diccionario de datos documenta el significado, el dominio y las unidades de cada elemento. Por ejemplo, una cantidad de productos no debe confundirse con un precio. En problemas monetarios conviene definir si se trabajará con céntimos enteros o decimales exactos y en qué momento se redondeará.

\begin{tablalibro}{3}{L{0.26}L{0.26}L{0.48}}
\caption{Diccionario de datos del pedido}\label{tab:diccionario-de-datos-del-pedido}\\
\toprule
\rowcolor{Papel}
\CabeceraTabla{Identificador} & \CabeceraTabla{Tipo conceptual} & \CabeceraTabla{Dominio o significado} \\
\midrule
\endfirsthead
\SigueTabla{3}
\toprule
\rowcolor{Papel}
\CabeceraTabla{Identificador} & \CabeceraTabla{Tipo conceptual} & \CabeceraTabla{Dominio o significado} \\
\midrule
\endhead
\midrule
\endfoot
\bottomrule
\multicolumn{3}{@{}p{\linewidth}@{}}{\NotaLibro{Elaboración propia.}}\\
\endlastfoot

cantidad & Entero & Mayor que cero \\
precio & Decimal & Soles, mayor o igual a cero \\
subtotal & Decimal & cantidad × precio \\
tasa & Decimal & Entre 0 y 1 \\
total & Decimal & subtotal × (1 \ensuremath{-} tasa) \\
\end{tablalibro}


\section{Diseño de estrategias para resolver problemas}

Diseñar una estrategia consiste en seleccionar un método de resolución compatible con los datos y las restricciones. La estrategia debe justificarse antes de convertirla en instrucciones detalladas. Los ejemplos siguientes constituyen una síntesis didáctica de técnicas elementales, no una clasificación exhaustiva de algoritmos.

\begin{tablalibro}{3}{L{0.28}L{0.34}L{0.38}}
\caption{Estrategias elementales de diseño}\label{tab:estrategias-elementales-de-diseno}\\
\toprule
\rowcolor{Papel}
\CabeceraTabla{Estrategia} & \CabeceraTabla{Aplicación} & \CabeceraTabla{Ejemplo} \\
\midrule
\endfirsthead
\SigueTabla{3}
\toprule
\rowcolor{Papel}
\CabeceraTabla{Estrategia} & \CabeceraTabla{Aplicación} & \CabeceraTabla{Ejemplo} \\
\midrule
\endhead
\midrule
\endfoot
\bottomrule
\multicolumn{3}{@{}p{\linewidth}@{}}{\NotaLibro{Elaboración propia.}}\\
\endlastfoot

Cálculo directo & Existe una relación matemática conocida. & Área de un rectángulo. \\
Selección por condiciones & La acción depende de una regla. & Clasificar una nota. \\
Recorrido secuencial & Se examina una colección de elementos. & Sumar gastos. \\
Búsqueda lineal & Se necesita localizar un valor en una lista sin ordenar. & Buscar un código de producto. \\
Acumulación y conteo & Se requiere un resumen de varios registros. & Calcular promedio de ventas. \\
\end{tablalibro}


\subsection{Comparación de alternativas}

Para localizar un código en una lista sin ordenar, una búsqueda lineal examina los elementos uno a uno y puede detenerse al encontrar el objetivo. Una búsqueda binaria reduce el intervalo a la mitad en cada paso, pero exige que los elementos estén ordenados y que sea posible acceder a su posición intermedia. Si solo se realizará una consulta, ordenar previamente podría añadir un costo innecesario; si se realizarán muchas, puede resultar conveniente.

La elección no depende únicamente del número de instrucciones escritas. Importan el volumen de datos, la frecuencia de ejecución, las precondiciones y la facilidad para verificar la solución. En los primeros diseños se priorizará la corrección y la claridad, sin perder de vista el costo de las operaciones.

\subsection{Tres estructuras básicas}

La secuencia ejecuta acciones en orden; la selección permite elegir una rama según una condición; la repetición permite ejecutar un bloque varias veces. \citet[pp.~86–95]{farrell2013} estudia estas estructuras como base de la programación estructurada. Combinarlas permite expresar soluciones complejas de forma organizada.

\section{Descomposición de soluciones y refinamiento paso a paso}

La descomposición de una solución divide el procedimiento general en tareas con responsabilidades delimitadas. El refinamiento sucesivo transforma cada tarea abstracta en acciones más concretas, hasta alcanzar instrucciones suficientemente precisas para su implementación. No debe confundirse con la descomposición del problema: aquella identifica necesidades; esta organiza las acciones para satisfacerlas.

La modularización facilita la comprensión y la reutilización de partes de la solución \citep[pp.~48–51]{farrell2013}. En un diseño inicial puede representarse cada módulo mediante su nombre, entradas y resultados, aun cuando todavía no se haya elegido el lenguaje de programación.

\subsection{Ejemplo: informe de asistencia}

Una institución desea calcular el porcentaje de asistencia de un estudiante. Se supone que cada registro válido corresponde a una sesión programada y contiene P (presente) o A (ausente). Una lista vacía se informa como ausencia de datos, no como 0 \% de asistencia.

\begin{tablalibro}{2}{L{0.31}L{0.69}}
\caption{Refinamiento de la solución de asistencia}\label{tab:refinamiento-de-la-solucion-de-asistencia}\\
\toprule
\rowcolor{Papel}
\CabeceraTabla{Nivel} & \CabeceraTabla{Acciones} \\
\midrule
\endfirsthead
\SigueTabla{2}
\toprule
\rowcolor{Papel}
\CabeceraTabla{Nivel} & \CabeceraTabla{Acciones} \\
\midrule
\endhead
\midrule
\endfoot
\bottomrule
\multicolumn{2}{@{}p{\linewidth}@{}}{\NotaLibro{Elaboración propia.}}\\
\endlastfoot

0 & Generar informe de asistencia. \\
1 & Leer registros; validar; calcular; presentar resultado. \\
2: validar & Comprobar que cada registro sea P o A y que la lista no esté vacía. \\
2: calcular & Inicializar presentes en cero; recorrer registros; incrementar si es P; dividir presentes entre total y multiplicar por 100. \\
2: presentar & Mostrar cantidad de sesiones, presentes y porcentaje. \\
\end{tablalibro}


\FiguraLibro{05-descomposicion-informe}{Descomposición jerárquica del informe de asistencia}{fig:descomposicion-informe}{Elaboración propia.}

Una gráfica de jerarquía muestra responsabilidades, no necesariamente el orden temporal de todas las instrucciones. Para especificar ese orden se emplea pseudocódigo o un diagrama de flujo. El refinamiento finaliza cuando las tareas pueden expresarse sin recurrir a instrucciones ambiguas como «procesar correctamente».

\section{Representación de algoritmos mediante pseudocódigo}

El pseudocódigo es una representación textual de la lógica de una solución que utiliza convenciones cercanas al lenguaje natural y a las instrucciones de programación. \citet[pp.~14–15]{farrell2013} explica su flexibilidad como herramienta de planificación. No existe una sintaxis universal; por ello, en esta obra se adoptan convenciones en español, sin exigir compatibilidad directa con un intérprete particular.

\begin{tablalibro}{2}{L{0.36}L{0.64}}
\caption{Convenciones de pseudocódigo del libro}\label{tab:convenciones-de-pseudocodigo-del-libro}\\
\toprule
\rowcolor{Papel}
\CabeceraTabla{Notación} & \CabeceraTabla{Significado} \\
\midrule
\endfirsthead
\SigueTabla{2}
\toprule
\rowcolor{Papel}
\CabeceraTabla{Notación} & \CabeceraTabla{Significado} \\
\midrule
\endhead
\midrule
\endfoot
\bottomrule
\endlastfoot

\Codigo{INICIO / FIN} & Delimitan el algoritmo. \\
\Codigo{LEER / MOSTRAR} & Reciben datos y presentan resultados. \\
\Codigo{\ensuremath{\leftarrow}} & Asigna un valor a una variable. \\
\Codigo{SI / SINO / FIN SI} & Representan una decisión. \\
\Codigo{PARA / FIN PARA} & Repiten un bloque para un rango definido. \\
\Codigo{MIENTRAS / FIN MIENTRAS} & Repiten mientras una condición sea verdadera. \\
\Codigo{Y, O, NO} & Operadores lógicos. \\
\Codigo{DIV, MOD} & Cociente entero y residuo. \\
\end{tablalibro}


\subsection{Estructura secuencial}

\Pseudocodigo{02-secuencia.txt}

Se presupone que base y altura son números positivos, expresados en la misma unidad. El resultado se expresa en unidades cuadradas. Las instrucciones se ejecutan una vez y en el orden indicado.

\subsection{Estructura selectiva}

\Pseudocodigo{02-seleccion.txt}

El saldo inicial se supone no negativo. La rama verdadera realiza el retiro y la falsa conserva el saldo. La igualdad entre importe y saldo debe admitirse, porque deja un saldo de cero.

\subsection{Estructura repetitiva}

\Pseudocodigo{02-repeticion.txt}

Se presupone que n es un entero no negativo y cada gasto es no negativo. El contador i toma valores de 1 a n, inclusive; cuando n es cero no hay iteraciones. El acumulador total conserva la suma de los valores leídos. La inicialización previa al ciclo evita utilizar un valor desconocido.

\section{Representación gráfica mediante diagramas de flujo}

El diagrama de flujo representa gráficamente las operaciones y el control de un algoritmo. \citet[pp.~16–19]{farrell2013} presenta símbolos para distinguir entrada, proceso, salida y dirección del flujo. Los diagramas ayudan a identificar rutas alternativas, conexiones incorrectas y ciclos sin condiciones de terminación.

\FiguraLibro{06-simbolos-flujo}{Símbolos elementales de un diagrama de flujo}{fig:simbolos-flujo}{Elaboración propia.}

El óvalo delimita el comienzo o final; el paralelogramo representa entrada o salida; el rectángulo indica una operación; el rombo formula una condición. Las flechas expresan el orden de ejecución. En una decisión binaria deben rotularse las ramas, por ejemplo, Sí y No. Un conector puede reducir cruces cuando el diagrama es extenso.

\subsection{Ejemplo: clasificación de temperatura}

Un sensor informa una temperatura numérica en grados Celsius. Como regla ilustrativa, se genera una alerta si la temperatura es mayor o igual a 30 °C; de lo contrario se informa un estado normal. No se trata de un umbral médico ni de seguridad industrial. La \RefFigura{fig:clasificacion-umbral} muestra las dos ramas de la decisión.

\FiguraLibro{07-clasificacion-umbral}{Diagrama de flujo de una clasificación por umbral}{fig:clasificacion-umbral}{Elaboración propia.}

Para comprobar la correspondencia con el pseudocódigo, cada símbolo debe tener una instrucción equivalente y cada rama debe conducir a un resultado. Las entradas 29, 30 y 31 permiten examinar ambos caminos y el límite de la condición.

\section{Seguimiento manual de algoritmos mediante tablas de prueba}

La prueba de escritorio consiste en simular manualmente las instrucciones del algoritmo, registrando los valores de las variables y las decisiones tomadas. \citet[p.~10]{farrell2013} la presenta como una forma de revisar la lógica antes de codificar. Una tabla de seguimiento documenta cada cambio relevante; no debe limitarse a mostrar el resultado final.

\subsection{Ejemplo: suma de tres gastos}

\Pseudocodigo{02-suma-gastos.txt}

\begin{tablalibro}{5}{L{0.26}R{0.08}R{0.15}R{0.15}L{0.36}}
\caption{Prueba de escritorio con gastos de S/ 12, 8 y 15}\label{tab:prueba-de-escritorio-con-gastos-de-s-12-8-y-15}\\
\toprule
\rowcolor{Papel}
\CabeceraTabla{Paso} & \CabeceraTabla{i} & \CabeceraTabla{gasto} & \CabeceraTabla{total} & \CabeceraTabla{Observación} \\
\midrule
\endfirsthead
\SigueTabla{5}
\toprule
\rowcolor{Papel}
\CabeceraTabla{Paso} & \CabeceraTabla{i} & \CabeceraTabla{gasto} & \CabeceraTabla{total} & \CabeceraTabla{Observación} \\
\midrule
\endhead
\midrule
\endfoot
\bottomrule
\endlastfoot

Inicialización & — & — & 0 & Antes del ciclo \\
Iteración 1 & 1 & 12 & 12 & 0 + 12 \\
Iteración 2 & 2 & 8 & 20 & 12 + 8 \\
Iteración 3 & 3 & 15 & 35 & 20 + 15 \\
Salida & — & — & 35 & Mostrar total \\
\end{tablalibro}


El resultado esperado, calculado de manera independiente, es S/ 35. Si el acumulador se inicializara dentro del ciclo, el resultado sería S/ 15, correspondiente al último gasto. El seguimiento revela el punto exacto en que la lógica se aparta del objetivo.

\subsection{Seguimiento de decisiones y ciclos}

En una selección conviene registrar si la condición evaluada es verdadera o falsa y la rama elegida. En una repetición debe verificarse la inicialización, la condición de continuidad, la actualización de las variables y la salida. Si un ciclo MIENTRAS no modifica ningún valor que permita volver falsa su condición, puede no terminar.

\begin{tablalibro}{5}{L{0.19}R{0.21}C{0.20}R{0.14}R{0.26}}
\caption{Seguimiento de un ciclo de cuenta regresiva}\label{tab:seguimiento-de-un-ciclo-de-cuenta-regresiva}\\
\toprule
\rowcolor{Papel}
\CabeceraTabla{Paso} & \CabeceraTabla{contador antes} & \CabeceraTabla{¿contador > 0?} & \CabeceraTabla{Salida} & \CabeceraTabla{contador después} \\
\midrule
\endfirsthead
\SigueTabla{5}
\toprule
\rowcolor{Papel}
\CabeceraTabla{Paso} & \CabeceraTabla{contador antes} & \CabeceraTabla{¿contador > 0?} & \CabeceraTabla{Salida} & \CabeceraTabla{contador después} \\
\midrule
\endhead
\midrule
\endfoot
\bottomrule
\multicolumn{5}{@{}p{\linewidth}@{}}{\NotaLibro{Elaboración propia.}}\\
\endlastfoot

1 & 3 & Sí & 3 & 2 \\
2 & 2 & Sí & 2 & 1 \\
3 & 1 & Sí & 1 & 0 \\
4 & 0 & No & — & 0 \\
\end{tablalibro}


La tabla supone que la instrucción de actualización resta una unidad al contador después de mostrarlo. Se observa que el ciclo finaliza al llegar a cero; si el valor inicial es cero, no se muestra ningún número.

\section{Validación de la lógica y comprobación de resultados}

La comprobación de un algoritmo busca determinar si sus instrucciones producen los resultados establecidos en la especificación. Se distingue entre verificar la coherencia interna del diseño y validar que la solución corresponda a la necesidad planteada. La prueba de escritorio y los casos de prueba aportan evidencia, pero un conjunto finito de ejemplos no demuestra por sí solo la corrección para todos los datos posibles.

\subsection{Diseño de casos de prueba}

Un caso de prueba contiene entradas, condiciones iniciales y un resultado esperado obtenido independientemente del algoritmo. Para evaluar una regla que acepta edades enteras de 18 a 65 años, se necesitan ejemplos normales, valores en los límites y entradas inválidas. La regla es ficticia y se utiliza exclusivamente para practicar condiciones.

\begin{tablalibro}{3}{L{0.20}L{0.26}L{0.54}}
\caption{Casos de prueba de una regla de intervalo}\label{tab:casos-de-prueba-de-una-regla-de-intervalo}\\
\toprule
\rowcolor{Papel}
\CabeceraTabla{Entrada} & \CabeceraTabla{Categoría} & \CabeceraTabla{Resultado esperado} \\
\midrule
\endfirsthead
\SigueTabla{3}
\toprule
\rowcolor{Papel}
\CabeceraTabla{Entrada} & \CabeceraTabla{Categoría} & \CabeceraTabla{Resultado esperado} \\
\midrule
\endhead
\midrule
\endfoot
\bottomrule
\multicolumn{3}{@{}p{\linewidth}@{}}{\NotaLibro{Elaboración propia.}}\\
\endlastfoot

35 & Normal válido & Aceptado \\
18 & Límite inferior & Aceptado \\
65 & Límite superior & Aceptado \\
17 & Fuera del intervalo & No aceptado \\
66 & Fuera del intervalo & No aceptado \\
-1 & Inválido & Error de dato \\
Texto & Inválido & Error de dato \\
\end{tablalibro}


Para este ejemplo, se considera inválida una edad negativa o no entera. Una edad no negativa fuera del intervalo es un dato válido que no satisface la regla de aceptación. Esta diferencia impide confundir los errores de entrada con los resultados legítimos de una decisión.

\subsection{Argumentos de corrección}

Una forma elemental de justificar un ciclo de suma consiste en observar que, después de procesar k elementos, el acumulador contiene la suma de los primeros k valores. Antes del ciclo, k es cero y la suma es cero. Cada iteración incorpora el siguiente valor, conservando esa propiedad. Cuando k alcanza la longitud de la lista, el acumulador contiene la suma completa. Si el contador avanza una posición cada vez, el ciclo termina. Esta propiedad que se conserva se denomina invariante.

La revisión final debe comprobar que todas las entradas requeridas están definidas, que las variables tienen valores antes de usarse, que las ramas cubren los casos previstos, que los ciclos terminan y que los resultados coinciden con los criterios de aceptación. Si se descubre un error, se modifica el diseño y se repiten las pruebas pertinentes.

\SeccionLibro[Caso práctico]{Caso práctico. Diseño del algoritmo de presupuesto}

\subsection{Control de presupuesto para una compra de materiales}

Se presenta una variante de implementación del caso del \RefCapitulo{cap:01-pensamiento-computacional}. Para mantener la continuidad conceptual, conserva el presupuesto de S/ 150 y el objetivo de decidir la viabilidad de la compra, pero utiliza céntimos enteros y una lista variable de productos. La variante permite concentrarse en el diseño algorítmico sin duplicar exactamente la tabla de materiales del capítulo anterior.

\begin{tablalibro}{2}{L{0.16}L{0.84}}
\caption{Requerimientos del caso práctico}\label{tab:requerimientos-del-caso-practico}\\
\toprule
\rowcolor{Papel}
\CabeceraTabla{Código} & \CabeceraTabla{Requerimiento} \\
\midrule
\endfirsthead
\SigueTabla{2}
\toprule
\rowcolor{Papel}
\CabeceraTabla{Código} & \CabeceraTabla{Requerimiento} \\
\midrule
\endhead
\midrule
\endfoot
\bottomrule
\multicolumn{2}{@{}p{\linewidth}@{}}{\NotaLibro{Elaboración propia.}}\\
\endlastfoot

R1 & Leer un presupuesto entero B no negativo y un número entero n de productos mayor o igual que cero. \\
R2 & Leer cantidad q entera positiva y precio unitario p entero no negativo en céntimos para cada producto. \\
R3 & Calcular cada subtotal y acumular el costo total. \\
R4 & Calcular saldo B \ensuremath{-} total e informar si la compra es viable (saldo \ensuremath{\geq} 0). \\
R5 & Ante un dato inválido, mostrar error y terminar sin emitir resultado de compra. \\
\end{tablalibro}


Se supone que los datos recibidos son numéricos enteros; en una implementación real será necesario comprobar además el formato antes de convertirlos. El presupuesto y el precio pueden ser cero. La cantidad de cada producto debe ser positiva, pero la lista vacía es válida y representa una compra sin artículos; en ese caso el total es cero.

\subsection{Modelo EPS y estrategia}

\begin{tablalibro}{3}{L{0.22}L{0.51}L{0.27}}
\caption{Modelo EPS del caso integrador}\label{tab:modelo-eps-del-caso-integrador}\\
\toprule
\rowcolor{Papel}
\CabeceraTabla{Entradas} & \CabeceraTabla{Procesamiento} & \CabeceraTabla{Salidas} \\
\midrule
\endfirsthead
\SigueTabla{3}
\toprule
\rowcolor{Papel}
\CabeceraTabla{Entradas} & \CabeceraTabla{Procesamiento} & \CabeceraTabla{Salidas} \\
\midrule
\endhead
\midrule
\endfoot
\bottomrule
\endlastfoot

\(B\), \(n\), registros (nombre, \(q\), \(p\)) & Validar datos, iniciar total en cero, sumar \(q \times p\) por material y comparar el total con \(B\). & Total, saldo y condición de viabilidad \\
Dato inválido & Detener el procedimiento & Mensaje de error \\
\end{tablalibro}


Se selecciona un recorrido secuencial con acumulación, pues cada producto contribuye una vez al total. La comparación con el presupuesto se efectúa después de procesar todos los productos válidos. No se necesita almacenar todos los subtotales para obtener el resultado final.

\subsection{Pseudocódigo de la solución}

\Pseudocodigo{02-control-presupuesto.txt}

TERMINAR indica la finalización inmediata del algoritmo y forma parte de las convenciones adoptadas en este libro. Si la lista está vacía, el algoritmo produce total cero y continúa con la evaluación del presupuesto. El saldo negativo expresa el importe faltante mediante su valor absoluto. En la salida al usuario se puede convertir cada importe de céntimos a soles, manteniendo dos decimales. La \RefFigura{fig:flujo-presupuesto} representa la misma lógica.

\subsection{Diagrama de flujo}

\FiguraLibro{08-flujo-presupuesto}{Flujo del control de presupuesto}{fig:flujo-presupuesto}{Elaboración propia.}

\subsection{Prueba de escritorio y validación}

Se dispone de S/ 150,00 (15 000 céntimos). Se compran 5 cuadernos a S/ 8,00, 10 lápices a S/ 2,00 y 7 cartulinas a S/ 5,00. El total esperado es S/ 95,00 y el saldo es S/ 55,00.

\begin{tablalibro}{5}{L{0.26}R{0.16}R{0.16}R{0.19}R{0.23}}
\caption{Seguimiento del caso integrador (céntimos)}\label{tab:seguimiento-del-caso-integrador-centimos}\\
\toprule
\rowcolor{Papel}
\CabeceraTabla{i} & \CabeceraTabla{q} & \CabeceraTabla{p} & \CabeceraTabla{q × p} & \CabeceraTabla{total} \\
\midrule
\endfirsthead
\SigueTabla{5}
\toprule
\rowcolor{Papel}
\CabeceraTabla{i} & \CabeceraTabla{q} & \CabeceraTabla{p} & \CabeceraTabla{q × p} & \CabeceraTabla{total} \\
\midrule
\endhead
\midrule
\ContinuaTabla{5}
\endfoot
\bottomrule
\multicolumn{5}{@{}p{\linewidth}@{}}{\NotaLibro{Elaboración propia.}}\\
\endlastfoot

Inicio & — & — & — & 0 \\
1 & 5 & 800 & 4000 & 4000 \\
2 & 10 & 200 & 2000 & 6000 \\
3 & 7 & 500 & 3500 & 9500 \\
\end{tablalibro}


El saldo es \(15\,000 - 9\,500 = 5\,500\) céntimos. La compra es viable porque el saldo es mayor o igual a cero.

\begin{tablalibro}{2}{L{0.46}L{0.54}}
\caption{Comprobación de escenarios adicionales}\label{tab:comprobacion-de-escenarios-adicionales}\\
\toprule
\rowcolor{Papel}
\CabeceraTabla{Presupuesto y productos} & \CabeceraTabla{Resultado esperado} \\
\midrule
\endfirsthead
\SigueTabla{2}
\toprule
\rowcolor{Papel}
\CabeceraTabla{Presupuesto y productos} & \CabeceraTabla{Resultado esperado} \\
\midrule
\endhead
\midrule
\endfoot
\bottomrule
\multicolumn{2}{@{}p{\linewidth}@{}}{\NotaLibro{Elaboración propia.}}\\
\endlastfoot

\(B=4650\), materiales de referencia & Total 4650, saldo 0, viable. \\
\(B=4000\), materiales de referencia & Total 4650, saldo \ensuremath{-}650, no viable. \\
\(B=0\), dos lápices a precio cero & Total 0, saldo 0, viable. \\
\(B=0\), lista de materiales vacía & Total 0, saldo 0, viable. \\
\(B=1000\), lápiz con \(q=0\) y \(p=100\) & Error de cantidad. \\
\(B=1000\), lápiz con \(q=2\) y \(p=-1\) & Error de precio. \\
\end{tablalibro}


La corrección del cálculo se justifica porque después de cada iteración el acumulador contiene el costo de todos los productos procesados. El contador avanza hasta n, de modo que el recorrido termina. La decisión final clasifica todos los saldos posibles en dos ramas mutuamente excluyentes. Los datos inválidos conducen a una salida de error antes de generar un resultado de compra.

\SeccionLibro{Cierre del capítulo}

\subsection{Resumen}

El diseño algorítmico transforma una especificación en un procedimiento ordenado, preciso y comprobable. El análisis de requerimientos establece las reglas y restricciones; el modelo EPS identifica los datos y las transformaciones; la selección de estrategias y el refinamiento organizan la solución. El pseudocódigo y los diagramas de flujo representan una misma lógica mediante formas diferentes. Las pruebas de escritorio permiten seguir los valores y decisiones, mientras que los casos de prueba y los argumentos de corrección ayudan a comprobar el cumplimiento de la especificación. Estas actividades preparan la implementación mediante un lenguaje de programación.

\subsection{Términos clave}

\begin{tablalibro}{2}{L{0.31}L{0.69}}
\caption{Glosario del capítulo}\label{tab:glosario-del-capitulo-2}\\
\toprule
\rowcolor{Papel}
\CabeceraTabla{Término} & \CabeceraTabla{Definición} \\
\midrule
\endfirsthead
\SigueTabla{2}
\toprule
\rowcolor{Papel}
\CabeceraTabla{Término} & \CabeceraTabla{Definición} \\
\midrule
\endhead
\midrule
\endfoot
\bottomrule
\endlastfoot

Algoritmo & Procedimiento organizado para resolver un problema. \\
Requerimiento & Necesidad o propiedad que debe satisfacer la solución. \\
Precondición & Condición exigida antes de una operación. \\
Poscondición & Condición garantizada al finalizar correctamente. \\
EPS & Modelo de entrada, proceso y salida. \\
Refinamiento & Desarrollo progresivo de tareas en instrucciones concretas. \\
Pseudocódigo & Representación textual de la lógica de un algoritmo. \\
Diagrama de flujo & Representación gráfica de operaciones y control. \\
Acumulador & Variable que conserva un resultado parcial. \\
Prueba de escritorio & Simulación manual de la ejecución. \\
Caso de prueba & Entrada y resultado esperado para una comprobación. \\
Invariante & Propiedad que se mantiene durante las iteraciones. \\
\end{tablalibro}


\subsection{Preguntas de comprensión}

\begin{enumerate}
\item ¿Qué diferencia existe entre un algoritmo y un programa?
\item ¿Por qué la finitud es una propiedad importante?
\item ¿Cómo se distingue un requerimiento de un criterio de aceptación?
\item ¿Qué función cumplen las precondiciones y poscondiciones?
\item ¿Qué información contiene un modelo EPS?
\item ¿Cuándo resulta apropiada una búsqueda lineal?
\item ¿Cómo se relacionan descomposición y refinamiento?
\item ¿Qué ventajas y limitaciones presenta el pseudocódigo?
\item ¿Cómo se representa una decisión en un diagrama de flujo?
\item ¿En qué se diferencian un caso de prueba y una tabla de seguimiento?
\item ¿Por qué deben evaluarse los valores límite?
\item ¿Por qué las pruebas exitosas no garantizan por sí solas la corrección universal?
\end{enumerate}

\subsection{Ejercicios de diseño algorítmico}

\begin{enumerate}
\item Conversión de tiempo. Diseñe un algoritmo que convierta segundos enteros no negativos a horas, minutos y segundos. Presente el modelo EPS y pruebe 0, 59, 60 y 3601.
\item Calificaciones. Calcule un promedio ponderado de tres notas con pesos de 20 \%, 30 \% y 50 \%. Defina el dominio, el criterio de aprobación y el pseudocódigo.
\item Inventario. Lea existencias y cantidad solicitada; determine si el pedido es posible y calcule el stock restante. Incluya igualdad y cantidades inválidas.
\item Consumo de agua. Registre los litros consumidos durante siete días y obtenga total y promedio. Represente el ciclo y construya una tabla de seguimiento.
\item Búsqueda. Dada una lista de códigos, encuentre la primera posición del código solicitado o indique que no existe. Pruebe una lista vacía y una con duplicados.
\item Tarifa de transporte. Diseñe una tarifa ficticia de S/ 3 por el primer kilómetro y S/ 1 por cada kilómetro entero adicional. Defina entradas válidas, pseudocódigo y pruebas de frontera.
\item Asistencia. Diseñe un algoritmo que calcule el porcentaje de presentes con registros P y A. Trate listas vacías y símbolos inválidos.
\end{enumerate}

\subsection{Desafíos de representación y seguimiento de algoritmos}

\textbf{Desafío 1.} Biblioteca escolar. Diseñe un algoritmo que determine si un préstamo está vencido a partir de dos fechas válidas, calcule los días de retraso y aplique una tarifa ficticia diaria. Declare las reglas, represente el flujo y pruebe un préstamo a tiempo, uno vencido y uno con fechas incoherentes.

\textbf{Desafío 2.} Monitoreo ambiental. Una estación registra n temperaturas. Diseñe una solución que calcule el máximo, el promedio y cuántas lecturas superan un umbral dado. Defina el tratamiento de una lista vacía y trace manualmente cuatro lecturas, incluyendo valores negativos y una lectura igual al umbral.

\textbf{Desafío 3.} Integración. Elija una necesidad identificada en el \RefCapitulo{cap:01-pensamiento-computacional}. Elabore su especificación, diccionario de datos, modelo EPS, refinamiento, pseudocódigo y diagrama de flujo. Prepare al menos cinco casos de prueba, incluyendo un límite y una entrada inválida. Explique por qué su algoritmo termina y cómo satisface los requerimientos.

Criterio de evaluación sugerido: claridad de la especificación (20 \%), coherencia del diseño (30 \%), correspondencia entre representaciones (20 \%), calidad de las pruebas (20 \%) y justificación de la solución (10 \%).

\subsection{Fuente de consulta del capítulo}

Las citas de este capítulo remiten a la edición de \citet{farrell2013}.

Los ejemplos, tablas, diagramas y actividades son de elaboración propia. Las obras de la bibliografía complementaria quedan pendientes de consulta directa.
